\documentclass[11pt]{article}
\usepackage[margin=1in]{geometry}
\usepackage{amsmath,amssymb,amsthm}
\usepackage{boldtensors}
\usepackage{hyperref}
\usepackage{booktabs}

\newcommand{\dd}{\mathrm{d}}

\title{Impedance-Matching vs. Robin-Robin \\
Coupling Conditions in Schwarz Methods}
\author{Alejandro Mota}

\begin{document}
\maketitle

\section{Setting}

Consider a solid body $\Omega \subset \mathbb{R}^3$ decomposed into two
subdomains $\Omega_1$ and $\Omega_2$ that are integrated independently
during a Schwarz iteration.  Two geometric configurations are used in
Norma.jl:

\paragraph{Non-overlap.}  The subdomains partition $\Omega$ and share an
interface $\Gamma = \partial\Omega_1 \cap \partial\Omega_2$ on which
coupling is imposed.  Let $~n_k$ be the outward unit normal to $\Omega_k$
on $\Gamma$ (so $~n_1 = -~n_2$), and let
\[
  ~t_k \;=\; ~\sigma_k\, ~n_k
\]
denote the traction on $\Omega_k$'s side of the interface.  Continuity of
the coupled solution requires
\begin{align}
  ~u_1 \;&=\; ~u_2 \qquad \text{(kinematic continuity),} \label{eq:continuity} \\
  ~t_1 + ~t_2 \;&=\; ~0 \qquad \text{(traction balance).} \label{eq:traction}
\end{align}
Classical Dirichlet-Neumann (DN) Schwarz picks up \eqref{eq:continuity}
on one side and \eqref{eq:traction} on the other.  The \emph{Robin-Robin}
(RR) and \emph{impedance} variants each replace both Dirichlet and Neumann
conditions with a linear combination of displacement and traction, which
offers extra tuning and, in the impedance case, wave absorption.

\paragraph{Overlap.}  The subdomains cover parts of $\Omega$ that
properly intersect: $\Omega_1 \cap \Omega_2 \ne \emptyset$.  Each $\Omega_k$
has a portion of its boundary,
$\Gamma_k \subset \Omega_j$ with $j \ne k$, that lies
strictly in the interior of the partner.  That boundary is not shared
with $\Omega_j$, so no mutual traction is available there; on the other
hand, the partner provides a well-defined interior state at every node of
$\Gamma_k$.  Classical overlap Schwarz enforces
kinematic matching
\[
  ~u_k \;=\; ~u_j \qquad \text{on } \Gamma_k,
\]
either pointwise (strong interpolation) or in an $L^2$ sense (weak
projection).  The \emph{impedance overlap} variant instead leaves the
overlap degrees of freedom free and drives $~u_k$ towards $~u_j$ through
an impedance-matching penalty, as detailed in Section~\ref{sec:overlap}.

\section{Robin-Robin (RR) Non-Overlap Schwarz}
\label{sec:rr-nonoverlap}

\subsection{Continuous formulation}

On each side of the interface, the user specifies a Robin parameter
$\alpha > 0$ and imposes
\begin{equation}
  ~t_k + \alpha\, ~u_k \;=\; ~g_k
    \qquad \text{on } \Gamma, \quad k = 1, 2, \label{eq:rr}
\end{equation}
with the forcing $~g_k$ set from the partner subdomain's state at the
previous Schwarz iterate:
\begin{equation}
  ~g_k^{(n+1)} \;=\; -\,~t_j^{(n)} + \alpha\, ~u_j^{(n)}
    \qquad j \ne k. \label{eq:rr-forcing}
\end{equation}
When the Schwarz iteration converges, $~u_k = ~u_j$ and
$~t_k = -~t_j$ so that \eqref{eq:rr} reproduces \eqref{eq:continuity}
and \eqref{eq:traction} simultaneously.  The combined
``elastic-spring'' interface smooths out the DN mismatch that drives DN
Schwarz towards a fixed point.

\subsection{Discrete implementation}

Let $~W = \int_\Gamma ~N\,~N^\top\,\dd\Gamma$ be the surface mass
(square-projector) matrix and let $~P$ be the Dirichlet projector
mapping the partner's interface trace to this subdomain's interface
nodes.  The RR boundary condition contributes to the solid subproblem on
$\Omega_k$ as follows:

\paragraph{Self stiffness term.} From \eqref{eq:rr}, the quantity
$\alpha\, ~u_k$ acts as an elastic boundary stiffness which is assembled
into the global tangent:
\begin{equation}
  ~K_{\mathrm{rs}} \;=\; \alpha\, ~W \qquad \text{(diagonal of }\Gamma).
\end{equation}

\paragraph{Partner contribution (right-hand side).}  The forcing
$~g_k$ from \eqref{eq:rr-forcing} is accumulated into the boundary-force
vector of $\Omega_k$:
\begin{equation}
  ~f_{\Gamma,k} \;\leftarrow\; ~f_{\Gamma,k}
     \;-\; ~t_j^{(n)}\big|_\Gamma
     \;+\; \alpha\, ~W\, ~P\, ~u_j^{(n)}.
\end{equation}
In Norma.jl the partner traction is negated once by
\texttt{get\_dst\_force} and the self term $\alpha\, ~W\, ~u_k$ is added
back via \texttt{apply\_robin\_bcs\_internal\_force!}.

\subsection{Parameter tuning}

$\alpha$ has units of stiffness per unit area
($\mathrm{Pa}/\mathrm{m} = \mathrm{N}/\mathrm{m}^3$).  It trades off
traction and displacement error:
\begin{itemize}
\item $\alpha \to \infty$: the Robin condition approaches a pure Dirichlet
      constraint; no traction feedback from the partner.
\item $\alpha \to 0$: the Robin condition approaches a pure Neumann
      (traction) constraint; no displacement feedback.
\item An optimal $\alpha$ minimizes the Schwarz convergence rate.
      Closed-form expressions have been studied in the optimized-Schwarz
      literature but depend on details of the discretization and
      problem.  In practice $\alpha$ is tuned empirically; the scaling
      $\alpha \sim E / h$ (Young's modulus over a representative element
      size) has been effective on the problems exercised here.
\end{itemize}

\subsection{Limitations for dynamic problems}

For dynamic (wave-propagation) problems, \eqref{eq:rr} behaves as a
purely elastic interface.  An incoming wave from $\Omega_j$ is partly
reflected back into $\Omega_j$ and partly transmitted into $\Omega_k$,
with reflection coefficient determined by the interface impedance
mismatch.  Because $\alpha$ has no physical correspondence to the
material's wave impedance, reflections in general are non-zero and can
drive energy growth when the two subdomains use different time
integrators (e.g.\ an implicit/explicit pairing across the interface).

\section{Impedance Non-Overlap Schwarz}
\label{sec:imp-nonoverlap}

\subsection{Motivation: non-reflecting boundary in 1D}

Consider a semi-infinite elastic bar with density $\rho$ and P-wave
modulus $M = \lambda + 2\mu$, supporting a right-going traveling wave
$u(x, t) = f(x - c_p t)$ with $c_p = \sqrt{M/\rho}$.  Differentiating
in $x$ and $t$ gives the kinematic identity $\dot u = u_t = -c_p\, u_x$.
Combined with 1D Hooke's law $\sigma = M\, u_x$ this yields
\begin{equation}
  \sigma \;=\; M\, u_x \;=\; -\, \tfrac{M}{c_p}\, \dot u
         \;=\; -\, \rho\, c_p\, \dot u.
\end{equation}
Define the \emph{characteristic acoustic impedance}
\begin{equation}
  Z \;=\; \rho\, c_p \;=\; \sqrt{\rho\,(\lambda + 2\mu)}.
\end{equation}
Then for a purely outgoing (right-going) wave,
\begin{equation}
  \sigma + Z\, \dot u \;=\; 0, \label{eq:absorbing-1d}
\end{equation}
which is the Lysmer-Kuhlemeyer non-reflecting boundary condition.
Imposing \eqref{eq:absorbing-1d} on the interface prevents spurious
reflections of outgoing waves and is the key ingredient in the
impedance variant.

\subsection{Continuous formulation}

Norma.jl's impedance Schwarz BC augments the Robin condition with the
impedance term from \eqref{eq:absorbing-1d}:
\begin{equation}
  ~t_k + Z\, \dot{~u}_k + \alpha\, ~u_k \;=\; ~g_k
  \qquad \text{on } \Gamma. \label{eq:imp}
\end{equation}
The forcing is built from the partner state:
\begin{equation}
  ~g_k^{(n+1)} \;=\; -\,~t_j^{(n)}
     \;+\; Z\, \dot{~u}_j^{(n)}
     \;+\; \alpha\, ~u_j^{(n)}.
\end{equation}
At convergence, $~u_k = ~u_j$, $\dot{~u}_k = \dot{~u}_j$, and
$~t_k = -~t_j$, so the impedance and displacement-penalty terms cancel
and \eqref{eq:imp} again reproduces continuity and traction balance.

Setting $\alpha = 0$ yields the pure impedance condition; the
displacement-penalty term $\alpha\, ~u$ is retained as an option to
regularize the quasi-static limit, where $\dot{~u} \to ~0$ and
\eqref{eq:imp} would otherwise reduce to the Neumann condition only.

\subsection{Discrete implementation}

For a Newmark integrator with parameters $\beta$, $\gamma$ and time step
$\Delta t$, the velocity is a linear function of displacement at the new
step,
\begin{equation}
  \dot{~u}_k^{(n+1)} \;=\; \frac{\gamma}{\beta\,\Delta t}\, ~u_k^{(n+1)}
     \;+\; (\text{terms known at } t_n). \label{eq:newmark-velo}
\end{equation}
Let $c_v \;=\; \gamma / (\beta\,\Delta t)$.  The self term
$Z\, \dot{~u}_k + \alpha\, ~u_k$ therefore produces the tangent
contribution
\begin{equation}
  ~K_{\mathrm{is}} \;=\; \bigl( Z\, c_v + \alpha \bigr)\, ~W,
     \label{eq:imp-stiffness}
\end{equation}
which is assembled into the global Hessian at each Newton iteration.
In the quasi-static limit Norma.jl sets $c_v = 0$, so
\eqref{eq:imp-stiffness} reduces to $\alpha\, ~W$ and the impedance BC
collapses onto RR.

The self internal-force contribution at the current state is
\begin{equation}
  ~f_{\mathrm{is}} \;=\; ~W\,\bigl(\, Z\, \dot{~u}_k \;+\; \alpha\, ~u_k \,\bigr),
\end{equation}
added by \texttt{build\_impedance\_schwarz\_force} (schwarz.jl:311-338).

The partner contribution is accumulated into the boundary-force vector:
\begin{equation}
  ~f_{\Gamma,k} \;\leftarrow\; ~f_{\Gamma,k}
     \;-\; ~t_j^{(n)}\big|_\Gamma
     \;+\; Z\, ~W\, ~P\, \dot{~u}_j^{(n)}
     \;+\; \alpha\, ~W\, ~P\, ~u_j^{(n)},
\end{equation}
as implemented in \texttt{apply\_bc\_detail} for
\texttt{SolidMechanicsImpedanceSchwarzBoundaryCondition}
(schwarz.jl:71-139).

\subsection{Derivation of $Z$ from material properties}

In Norma.jl the impedance is computed per-domain from the material
parameters:
\begin{equation}
  Z \;=\; \sqrt{\rho\, (\lambda + 2\mu)} \;=\; \rho\, c_p, \qquad
  \text{with } c_p = \sqrt{(\lambda + 2\mu)/\rho}.
\end{equation}
No user tuning is required for the impedance term; only the optional
displacement-penalty $\alpha$ is exposed via the \texttt{robin parameter}
field.

\section{Overlap Schwarz}
\label{sec:overlap}

On an overlap boundary $\Gamma_k$, the partner
field $~u_j$ is available but the partner \emph{traction} is not,
because $\Gamma_k$ is not a shared interface:
it is a cut through the interior of $\Omega_j$.  Only those coupling
conditions that do not require a partner traction can be imposed on an
overlap boundary.  Norma.jl provides two such families.

\subsection{Classical overlap (Dirichlet coupling)}

The canonical overlap Schwarz BC replaces the free boundary of $\Omega_k$
on $\Gamma_k$ with an imposed kinematic trace
drawn from the partner,
\begin{equation}
  ~u_k \;=\; ~u_j,
  \qquad
  \dot{~u}_k \;=\; \dot{~u}_j,
  \qquad
  \ddot{~u}_k \;=\; \ddot{~u}_j
  \qquad \text{on } \Gamma_k. \label{eq:overlap-dbc}
\end{equation}
The trace is constructed in one of two ways:

\paragraph{Strong interpolation.}  For each overlap node
$~x_i \in \Gamma_k$, find the partner element
containing $~x_i$ in its reference configuration, evaluate the partner
shape functions $~N_j(~x_i)$, and assign
\[
  ~u_k(~x_i) \;=\; \sum_a ~N_j^a(~x_i)\, ~u_j^a.
\]
The corresponding degrees of freedom on $\Omega_k$ are then frozen.

\paragraph{Weak ($L^2$) projection.}  A Dirichlet projector
$~P$ is precomputed by solving the boundary $L^2$ problem
$~W_k\, \hat{~u}_k = ~C_{kj}\, ~u_j$, where $~W_k$ is the mass matrix of
$\Gamma_k$ and $~C_{kj}$ is the cross-mass coupling
the partner's shape functions to the overlap boundary.  Overlap-node
displacements are then assigned $~u_k = ~P\, ~u_j$.

Both variants are implemented by
\texttt{SolidMechanicsOverlapSchwarzBoundaryCondition} and selected via
the \texttt{use\_weak} flag (\texttt{coupling\_strong\_dbc},
\texttt{coupling\_weak\_overlap\_dbc} in
\texttt{schwarz.jl:453--506}).  Because \eqref{eq:overlap-dbc} constrains
the overlap DOFs, classical overlap adds no interface tangent term: the
constrained rows and columns are eliminated from the system.

\subsection{Impedance overlap}

The impedance overlap BC, by contrast, leaves the overlap DOFs
\emph{free} and couples $\Omega_k$ to the partner through an
impedance/Robin penalty,
\begin{equation}
  Z\, \bigl(\dot{~u}_k - \dot{~u}_j\bigr)
  \;+\; \alpha\, \bigl(~u_k - ~u_j\bigr)
  \;=\; ~0
  \qquad \text{on } \Gamma_k. \label{eq:imp-overlap}
\end{equation}
No traction term appears on either side of \eqref{eq:imp-overlap},
which is the essential discrete difference from the non-overlap
impedance condition \eqref{eq:imp}.  On $\Omega_k$'s side, the elastic
equilibrium on $\Gamma_k$ supplies the natural
(free-surface) condition $~t_k = ~0$; \eqref{eq:imp-overlap} then
drives the self state towards the partner's state via a spring-dashpot
coupling tuned to the material impedance.

The discrete assembly is the same as in the non-overlap impedance
case, minus the partner traction:
\begin{itemize}
\item \textbf{Self tangent.}
      $~K_{\mathrm{is}} = (Z\, c_v + \alpha)\, ~W$, with
      $c_v = \gamma/(\beta\Delta t)$ under Newmark, assembled by
      \texttt{build\_impedance\_overlap\_schwarz\_stiffness}
      (schwarz.jl:265--293).
\item \textbf{Self internal force.}
      $~f_{\mathrm{is}} = ~W\,\bigl(Z\, \dot{~u}_k + \alpha\, ~u_k\bigr)$,
      assembled by \texttt{build\_impedance\_schwarz\_force}
      (schwarz.jl:311--338), shared with the non-overlap impedance BC.
\item \textbf{Partner right-hand side.}
      $~f_\Gamma \leftarrow ~f_\Gamma + ~W\,\bigl(Z\, \dot{~u}_j
      + \alpha\, ~u_j\bigr)$, where $\dot{~u}_j$ and $~u_j$ are obtained
      by strong interpolation from the partner element containing each
      overlap node (\texttt{apply\_bc\_detail} for
      \texttt{SolidMechanicsImpedanceOverlapSchwarzBoundaryCondition},
      schwarz.jl:216--263).
\end{itemize}
At the Schwarz fixed point $~u_k = ~u_j$ and $\dot{~u}_k = \dot{~u}_j$
on $\Gamma_k$; the self and partner impedance
contributions cancel, leaving the overlap surface as a
free (traction-free) surface whose interior solution equals the
partner's.  Because the impedance matches the material's own P-wave
impedance, outgoing waves in $\Omega_k$ see no reflection at
$\Gamma_k$: they exit $\Omega_k$ consistently
with the solution that continues inside $\Omega_j$.

Norma.jl additionally allows a step-dependent multiplier on $Z$ through
the \texttt{impedance\_scale} vector (schwarz.jl:206--214).  This is
useful for ramping the impedance coupling in during the early part of a
simulation to avoid exciting transients.

\subsection{Why there is no Robin-Robin overlap variant}

An RR condition, $~t + \alpha\, ~u = ~g$, requires a partner traction to
build its right-hand side $~g = -~t_j + \alpha\, ~u_j$.  On an overlap
boundary there is no partner traction to borrow.  Formally dropping the
$-~t_j$ term from RR yields
\[
  ~t_k + \alpha\, ~u_k \;=\; \alpha\, ~u_j
  \quad\Longleftrightarrow\quad
  ~t_k + \alpha\, (~u_k - ~u_j) \;=\; ~0,
\]
which is the $Z = 0$ special case of \eqref{eq:imp-overlap}: a pure
displacement penalty with no dynamic absorption.  This is subsumed by
the impedance overlap BC (set $Z = 0$, keep $\alpha$) and is not exposed
as a separate option.

\section{Comparison}

\subsection{Side-by-side formulation}

\paragraph{Non-overlap variants.}
\begin{center}
\small
\begin{tabular}{lcc}
\toprule
& RR & Impedance \\
\midrule
Boundary condition &
  $~t_k + \alpha\, ~u_k = ~g_k$ &
  $~t_k + Z\, \dot{~u}_k + \alpha\, ~u_k = ~g_k$ \\[0.3em]
Self tangent &
  $\alpha\, ~W$ &
  $(Z\, c_v + \alpha)\, ~W$ \\[0.3em]
Self internal force &
  $\alpha\, ~W\, ~u_k$ &
  $~W\,\bigl(Z\, \dot{~u}_k + \alpha\, ~u_k\bigr)$ \\[0.3em]
Partner RHS &
  $-~t_j + \alpha\, ~W\, ~P\, ~u_j$ &
  $-~t_j + Z\, ~W\, ~P\, \dot{~u}_j + \alpha\, ~W\, ~P\, ~u_j$ \\[0.3em]
Partner traction $~t_j$ &
  used &
  used \\[0.3em]
$Z$ set by &
  (not used) &
  material (auto) \\[0.3em]
$\alpha$ set by &
  user &
  user, optional \\[0.3em]
Quasi-static limit &
  same &
  collapses to RR \\
\bottomrule
\end{tabular}
\end{center}

\paragraph{Overlap variants.}
\begin{center}
\small
\begin{tabular}{lcc}
\toprule
& Classical (DBC) & Impedance \\
\midrule
Boundary condition &
  $~u_k = ~u_j$ (strong or $L^2$) &
  $Z\,(\dot{~u}_k - \dot{~u}_j) + \alpha\,(~u_k - ~u_j) = ~0$ \\[0.3em]
Self tangent &
  (DOFs eliminated) &
  $(Z\, c_v + \alpha)\, ~W$ \\[0.3em]
Self internal force &
  (none) &
  $~W\,\bigl(Z\, \dot{~u}_k + \alpha\, ~u_k\bigr)$ \\[0.3em]
Partner RHS &
  $~u_k \leftarrow ~P\, ~u_j$ &
  $~W\,\bigl(Z\, \dot{~u}_j + \alpha\, ~u_j\bigr)$ \\[0.3em]
Partner traction $~t_j$ &
  (not available) &
  (not used) \\[0.3em]
$Z$ set by &
  (not used) &
  material (auto), optional scale \\[0.3em]
$\alpha$ set by &
  (not used) &
  user, optional \\[0.3em]
Quasi-static limit &
  same &
  pure $\alpha$-penalty ($c_v = 0$) \\
\bottomrule
\end{tabular}
\end{center}

\subsection{Reflection at a non-overlap interface}

The reflection analysis below is for the non-overlap geometry, where
a single shared interface carries both the self and the partner
traction.  In the overlap case each subdomain's coupling boundary is
embedded in the partner's interior, so the analysis below does not
apply directly; Section~\ref{sec:overlap} discusses the overlap behavior
separately.

Place the interface at $x = 0$ and the subdomain $\Omega_k$ in $x < 0$,
so the outward normal of $\Omega_k$ at $\Gamma$ is $+\hat x$.  Look for
harmonic plane-wave solutions of the 1D wave equation: a right-going
incident wave $u_{\mathrm{inc}}$ propagating out of $\Omega_k$ toward
$\Gamma$, and a left-going reflected wave $u_{\mathrm{ref}}$
propagating back into $\Omega_k$,
\[
  u_{\mathrm{inc}}(x, t) = F\, \mathrm{e}^{\mathrm{i}(\omega t - k x)},
  \qquad
  u_{\mathrm{ref}}(x, t) = R\, \mathrm{e}^{\mathrm{i}(\omega t + k x)}.
\]
Here $\omega$ is the angular frequency, $k = \omega / c_p$ is the
wavenumber from the 1D elastic dispersion relation, and $F$, $R$ are
complex amplitudes.  Substituting into $\sigma = M\, u_x$ and using
$M\, k = M\, \omega / c_p = \rho c_p^2 \cdot \omega / c_p = Z\, \omega$
gives the interface values
\[
  u(0, t) = (F + R)\, \mathrm{e}^{\mathrm{i}\omega t},
  \quad
  \sigma(0, t) = \mathrm{i}\, \omega\, Z\, (R - F)\, \mathrm{e}^{\mathrm{i}\omega t},
  \quad
  \dot u(0, t) = \mathrm{i}\, \omega\, (F + R)\, \mathrm{e}^{\mathrm{i}\omega t}.
\]

\paragraph{RR boundary condition.}  Imposing $\sigma + \alpha\, u = 0$
at $x = 0$ gives
\[
  \mathrm{i}\, Z\, \omega\, (R - F) \;+\; \alpha\, (F + R) \;=\; 0,
\]
and solving for the reflection coefficient $R_{\mathrm{RR}}(\omega)
\equiv R / F$ yields
\begin{equation}
  R_{\mathrm{RR}}(\omega) \;=\;
     \frac{\mathrm{i}\, \omega\, Z - \alpha}
          {\mathrm{i}\, \omega\, Z + \alpha}.
  \label{eq:R-rr}
\end{equation}
For any real $\alpha > 0$ the numerator and denominator have the same
modulus, so $\lvert R_{\mathrm{RR}}(\omega) \rvert \equiv 1$: the
elastic interface \emph{fully reflects} the incoming wave with a
frequency-dependent phase shift.  No wave energy crosses the interface.

\paragraph{Impedance boundary condition.}  Imposing
$\sigma + Z\, \dot u + \alpha\, u = 0$ at $x = 0$ gives
\[
  \mathrm{i}\, Z\, \omega\, (R - F) \;+\; \mathrm{i}\, Z\, \omega\, (F + R)
  \;+\; \alpha\, (F + R) \;=\; 0,
\]
which simplifies to
$(2\, \mathrm{i}\, Z\, \omega + \alpha)\, R + \alpha\, F = 0$, so
\begin{equation}
  R_{\mathrm{imp}}(\omega) \;=\;
     -\, \frac{\alpha}{2\, \mathrm{i}\, \omega\, Z + \alpha}.
  \label{eq:R-imp}
\end{equation}
With $\alpha = 0$ the pure impedance condition is satisfied identically
by an outgoing wave, and $R_{\mathrm{imp}} \equiv 0$: the wave is
absorbed at all frequencies.  For $\alpha > 0$,
$\lvert R_{\mathrm{imp}}(\omega) \rvert
   = \alpha / \sqrt{\alpha^2 + 4\, \omega^2\, Z^2}$, which is strictly
below one at every frequency and tends to zero as $\alpha \to 0$.

\paragraph{Comparison.}  The contrast between \eqref{eq:R-rr} and
\eqref{eq:R-imp} is the key difference between the two Schwarz
variants.  Classical RR conserves energy at the interface
($\lvert R_{\mathrm{RR}} \rvert \equiv 1$), reflecting outgoing waves
back into the originating subdomain; impedance Schwarz dissipates
outgoing-wave energy into the partner
($\lvert R_{\mathrm{imp}} \rvert < 1$, and
$R_{\mathrm{imp}} = 0$ at $\alpha = 0$).  This is the mechanism by
which impedance Schwarz eliminates the reflection-driven instability
that plagues classical RR in time-dependent problems.

\subsection{Energy behavior with mixed integrators}

With classical RR coupling across an implicit/explicit interface,
reflections from the elastic interface can beat against each other and
cause non-monotone energy growth over many Schwarz iterations.  The
impedance variant shunts the outgoing wave energy into the partner and
removes that pathway, which stabilizes dynamic simulations that
integrate the two subdomains differently (e.g.\ Newmark + central
difference) without requiring a global time step.

\subsection{Quasi-static limit}

For quasi-static problems ($\dot{~u} \to ~0$, $c_v = 0$ in
\eqref{eq:imp-stiffness}), the impedance condition \eqref{eq:imp}
reduces exactly to the Robin condition \eqref{eq:rr}.  Norma.jl's
quasi-static branch computes $c_v = 0$ explicitly and the impedance
pipeline assembles the same $\alpha\, ~W$ tangent as RR.  In particular,
an impedance run with $\alpha = 0$ in the quasi-static limit degenerates
to a pure Neumann condition and does not couple the subdomains through
displacement, which is an important edge case to avoid in static use.

\section{Summary}

\begin{itemize}
\item On a \emph{non-overlap} interface, RR Schwarz enforces an elastic
      condition $~t + \alpha\, ~u = ~g$; $\alpha$ is a user-supplied
      stiffness that trades off traction and displacement error.
\item Impedance non-overlap Schwarz augments the Robin condition with the
      Lysmer-Kuhlemeyer absorbing term $Z\, \dot{~u}$ using the
      material's own P-wave impedance
      $Z = \sqrt{\rho(\lambda + 2\mu)}$.  No tuning is required for $Z$;
      the displacement-penalty $\alpha$ is optional.  Discretely, the
      only differences from RR are (i) an extra $Z\, c_v$ contribution
      to the interface tangent and (ii) a partner velocity term in the
      right-hand side; both are assembled with the same surface operator
      $~W$ and projection $~P$ used by RR.
\item In the quasi-static limit $c_v \to 0$, impedance non-overlap
      Schwarz collapses onto RR.  For dynamic problems it is
      non-reflecting at the interface
      ($\lvert R_{\mathrm{imp}} \rvert < 1$, with
      $R_{\mathrm{imp}} = 0$ at $\alpha = 0$), whereas RR is fully
      reflecting ($\lvert R_{\mathrm{RR}} \rvert \equiv 1$).  This is
      the mechanism that removes the reflection-driven instability of
      classical RR with mixed integrators.
\item On an \emph{overlap} boundary, no partner traction is available,
      so RR has no natural analogue.  Norma.jl provides two overlap
      variants: classical DBC coupling, which interpolates the partner
      displacement (strongly or in $L^2$) and constrains the overlap
      DOFs; and impedance overlap, which leaves the overlap DOFs free
      and couples to the partner through the spring-dashpot penalty
      $Z(\dot{~u}_k - \dot{~u}_j) + \alpha(~u_k - ~u_j) = ~0$.
\item The impedance overlap BC uses the same self-side tangent
      $(Z\, c_v + \alpha)\, ~W$ and self-side internal force
      $~W\,(Z\, \dot{~u}_k + \alpha\, ~u_k)$ as the non-overlap variant;
      the only discrete difference is that the partner right-hand side
      carries no $-~t_j$ term, because the overlap boundary is a cut
      through $\Omega_j$'s interior rather than a shared interface.  At
      convergence, outgoing waves in $\Omega_k$ exit through the overlap
      surface without reflection, matched to the solution that continues
      inside $\Omega_j$.
\end{itemize}

\end{document}
